% GENERATED LOCAL-BUILD FLAT PART — DO NOT MERGE

% ---- BEGIN TIR/monograph/v12/chapters/ch12_charged_leptons_neutrinos.tex ----
% !TEX root = ../../tir_monograph_v12.tex
% V12_SOURCE_MAP:
% - TIR/monograph/chapters/ch05_electron_action.tex
% - TIR/monograph/chapters/ch06_generation_release.tex
% - TIR/monograph/chapters/ch07_lepton_spectrum.tex
% - TIR/monograph/chapters/ch17_neutrino_masses.tex
% - TIR/monograph/chapters/ch30_parameter_summary.tex
% - TIR/monograph/chapters/ch32_pdg2026_validation_addendum.tex
% - TIR/foundations/TIR_COEFFICIENT_MAGNITUDE_PARENT_EVALUATION_V0_1.md
% Validator: TIR/validation/tir_v12_neutrino_action_consistency_v0_1.py
\chapter{Charged Leptons and Neutrinos}
\label{ch:v12_leptons_neutrinos}

\paragraph{Established context.}
Charged-lepton masses and neutrino oscillation observables are externally
measured quantities. The TIR relations collected here are sector constructions
built over the upstream normalization and coefficient architecture. Current
external-data comparison is centralized in Chapter~\ref{ch:v12_evidence_matrix}
rather than repeated in this chapter.

\section{Exponential mass map}
For the mass sectors considered here, TIR uses the dimensionless action map
\begin{equation}
\boxed{
m=E_P\exp\!\left(-\frac{S}{\kappa}\right),
}
\label{eq:v12-exponential-mass-map}
\end{equation}
where $E_P$ is an explicit external scale and $\kappa$ is the structural
normalization of Chapter~\ref{ch:v12_kappa_normalization}.

\vTwelveStatus{B}{--}{OPEN}
The parent-count arithmetic for the three historical charged-lepton coefficient
packets is now closed exactly in
Chapter~\ref{ch:v12_action_coefficients}. What remains open upstream is the
coefficient-free transition-parent selector, followed by the physical binding
and empirical test of the action-to-mass map. A parent-evaluation PASS does not
change the charged-lepton precision verdict below.

\section{Charged-lepton action family}
The historical electron action is
\begin{equation}
\boxed{
S_e
=\frac12-3\kappa+\frac{\kappa}{L_3}-\frac{\kappa^2}{2}.
}
\label{eq:v12-electron-action}
\end{equation}
The two generation-release operators are
\begin{align}
R_{e\mu}
&=5\kappa+\frac{2\kappa}{L_3}+\frac{L_3+1}{2}\kappa^2,
\label{eq:v12-electron-muon-release}\\
R_{\mu\tau}
&=3\kappa-\frac{\kappa}{L_3}-\frac{L_3}{2}\kappa^2.
\label{eq:v12-muon-tau-release}
\end{align}
The derived action chain is
\begin{equation}
\boxed{
S_\mu=S_e-R_{e\mu},
\qquad
S_\tau=S_\mu-R_{\mu\tau}.
}
\label{eq:v12-charged-lepton-actions}
\end{equation}
Using the canonical v12 values of $\kappa$ and $L_3=7$ gives
\begin{align}
S_e&=0.473691600121\ldots,\\
S_\mu&=0.424761179774\ldots,\\
S_\tau&=0.398790835923\ldots .
\end{align}
Inserted into Eq.~\eqref{eq:v12-exponential-mass-map}, the historical formula
family evaluates to
\begin{align}
m_e^{\rm TIR}&=0.511641997\ldots\ \mev,\\
m_\mu^{\rm TIR}&=104.831758\ldots\ \mev,\\
m_\tau^{\rm TIR}&=1767.503883\ldots\ \mev.
\end{align}
These numbers identify the exact formula version being discussed; the later
v11 summary contains a different retrospectively refined numerical triplet.
The current publication comparison therefore lives only in the unified evidence
matrix.

\vTwelveStatus{C}{RETROSPECTIVE}{FAIL}
The 2026 precision audit classifies the charged-lepton sector as arithmetic
coherent but not precision-level agreement with experiment.

\section{Neutrino tetrahedral action branch}
The historical neutrino branch starts from
\begin{equation}
S_{\rm bare}
=\frac{1+L_4/L_3}{2}
=\frac9{14},
\label{eq:v12-neutrino-bare-action}
\end{equation}
and the projected tetrahedral face factor
\begin{equation}
A_{\rm face}
=\frac{(L_4/L_3)^2}{2}
=\frac2{49}.
\label{eq:v12-neutrino-face-area}
\end{equation}
Define the action offset
\begin{equation}
\boxed{
dS
=\kappa A_{\rm face}(1-\kappa).
}
\label{eq:v12-neutrino-offset}
\end{equation}

\section{Formula--value reconciliation}
The v11 neutrino chapter prints
\[
S_1=S_{\rm bare}+\kappa dS,
\]
but its published mass values are generated by the single-offset relation
\begin{equation}
\boxed{
S_1^{\rm rec}=S_{\rm bare}+dS.
}
\label{eq:v12-neutrino-reconciled-S1}
\end{equation}
The deterministic diagnostic
\path{TIR/validation/tir_v12_neutrino_action_consistency_v0_1.py}
checks both equations. Literal evaluation of the printed double-$\kappa$ rule
moves all three printed masses by approximately $4.088\%$, whereas
Eq.~\eqref{eq:v12-neutrino-reconciled-S1} reproduces the historical mass triplet
to the displayed rounding precision.

\vTwelveStatus{D}{RETROSPECTIVE}{PASS}
This is a formula/value consistency diagnosis and reconstruction of the
historical calculation path. It does not promote the reconstructed relation to
an upstream structural theorem.

\section{Neutrino ratio pattern}
The declared mass-ratio labels are
\begin{equation}
\boxed{
r_i\in\{1,L_4,L_3+L_4+1\}=\{1,2,10\}.
}
\label{eq:v12-neutrino-ratios}
\end{equation}
With the reconciled historical action base,
\begin{equation}
S_i=S_1^{\rm rec}-\kappa\ln r_i,
\label{eq:v12-neutrino-actions}
\end{equation}
so the exponential map gives exactly the ratio inheritance
\begin{equation}
\frac{m_i}{m_1}=r_i.
\label{eq:v12-neutrino-ratio-inheritance}
\end{equation}
Numerically this historical formula version gives
\begin{align}
m_1&=0.005009994\ldots\ \eV,\\
m_2&=0.010019988\ldots\ \eV,\\
m_3&=0.050099939\ldots\ \eV,
\end{align}
and therefore
\begin{align}
\Delta m_{21}^2
&=7.5300117\times10^{-5}\ \eV^2,\\
\Delta m_{31}^2
&=2.4849039\times10^{-3}\ \eV^2.
\end{align}

\vTwelveStatus{C}{RETROSPECTIVE}{COMPATIBLE}
The mass-splitting comparison remains a retrospective compatibility surface.
Absolute mass values remain a model output whose external test status is tracked
separately in Chapter~\ref{ch:v12_evidence_matrix}.

\section{Separation from flavour mixing}
PMNS angles and CP phase are not folded into the present mass-action chapter.
They are owned by Chapter~\ref{ch:v12_ckm_pmns}, where the three-flavour carrier
of Chapter~\ref{ch:v12_three_flavour_su3} is already available upstream. This
keeps neutrino mass-scale construction distinct from the flavour-mixing
observable matrix.

% ---- END TIR/monograph/v12/chapters/ch12_charged_leptons_neutrinos.tex ----
\clearpage

% ---- BEGIN TIR/monograph/v12/chapters/ch13_ckm_pmns_flavour_mixing.tex ----
% !TEX root = ../../tir_monograph_v12.tex
% V12_SOURCE_MAP:
% - TIR/monograph/chapters/ch16_pmns_mixing.tex
% - TIR/monograph/chapters/ch18_ckm_matrix.tex
% - TIR/monograph/chapters/ch19_ckm_cp_violation.tex
% - TIR/validation/pdg2026/TIR_PDG2026_VALIDATION_MATRIX_V2.md
% Validators:
% - TIR/validation/tir_v12_flavour_mixing_internal_consistency_v0_1.py
% - TIR/validation/tir_ckm_jarlskog_invariant_closure_v0_1.py
\chapter{CKM/PMNS and Flavour Mixing}
\label{ch:v12_ckm_pmns}

\paragraph{Established context.}
Lepton mixing is conventionally described by the PMNS framework and quark
mixing by the CKM framework.  This chapter collects the TIR arithmetic/geometric
assignments only after the common three-flavour carrier and its $SU(3)_F$
structure have been established in Chapter~\ref{ch:v12_three_flavour_su3}.
External-data comparison is owned by Chapter~\ref{ch:v12_evidence_matrix}.

\section{PMNS structural assignments}
The historical TIR PMNS map uses the discrete labels of
Chapter~\ref{ch:v12_discrete_labels}.  The retained formula family is
\begin{align}
\sin\theta_{13}
&=\frac1{L_3}=\frac17,
\label{eq:v12-pmns-theta13}\\
\sin^2\theta_{12}
&=\frac{L_4}{L_3+L_4}
+\left(\frac{L_4}{L_3}\right)^2,
\label{eq:v12-pmns-theta12}\\
\sin^2\theta_{23}
&=\frac12+\frac{L_4}{L_3^2},
\label{eq:v12-pmns-theta23}\\
\delta_{\rm CP}^{\rm PMNS}
&=\pi\left(
1+\frac{L_4}{L_3}
+\left(\frac{L_4}{L_3}\right)^2
\right).
\label{eq:v12-pmns-delta}
\end{align}
For $(L_3,L_4)=(7,2)$ these evaluate to
\begin{align}
\sin^2\theta_{12}&=0.3038548753\ldots,\\
\sin^2\theta_{23}&=0.5408163265\ldots,\\
\sin^2\theta_{13}&=\frac1{49}=0.0204081633\ldots,\\
\delta_{\rm CP}^{\rm PMNS}&=246.122449^\circ\ldots .
\end{align}

\subsection{Current PMNS verdicts}
The PDG-2026 validation matrix assigns different evidential states to the four
observables; v12 therefore does not compress them into one sector-wide claim.

{\setlength{\tabcolsep}{3pt}
\begin{longtable}{p{0.32\textwidth}p{0.13\textwidth}p{0.22\textwidth}p{0.20\textwidth}}
\toprule
Observable & Claim Class & Timing & Verdict\\
\midrule
$\sin^2\theta_{12}$ & C & RETROSPECTIVE & COMPATIBLE\\
$\sin^2\theta_{23}$ & C & RETROSPECTIVE & OPEN\\
$\sin^2\theta_{13}=1/49$ & C & RETROSPECTIVE & TENSION\\
$\delta_{\rm CP}^{\rm PMNS}$ & C & RETROSPECTIVE & OPEN\\
\bottomrule
\end{longtable}}
The reactor-angle tension is retained without retuning the formula.

\section{CKM primary assignments}
The TIR Cabibbo parameter is
\begin{equation}
\boxed{
\lambda
=\frac{L_4}{L_3+L_4}
+\frac{L_4}{L_3}\kappa
=0.2248488365\ldots .
}
\label{eq:v12-ckm-lambda}
\end{equation}
The two small primary magnitudes used downstream are
\begin{align}
|V_{cb}|
&=\frac{(L_4/L_3)^2}{2}
=\frac2{49}
=0.0408163265\ldots,
\label{eq:v12-ckm-vcb}\\
|V_{ub}|
&=\frac{(L_4/L_3)^2L_4}{(L_3+L_4)L_5}
=\frac8{2205}
=0.0036281179\ldots .
\label{eq:v12-ckm-vub}
\end{align}
The corresponding Wolfenstein coefficient is
\begin{equation}
A=\frac{|V_{cb}|}{\lambda^2}
=0.8073328542\ldots .
\label{eq:v12-ckm-A}
\end{equation}

The displayed historical CKM magnitude matrix is
\begin{equation}
|V_{\rm CKM}|_{\rm hist}
=\begin{pmatrix}
0.97439 & 0.22485 & 0.00363\\
0.22470 & 0.97357 & 0.04082\\
0.00886 & 0.04001 & 0.99916
\end{pmatrix}.
\label{eq:v12-ckm-magnitude-matrix}
\end{equation}
At the printed precision its row- and column-norm residuals are below
$7\times10^{-6}$.  Full complex unitarity and phase-convention statements are
kept distinct from this rounded magnitude check.

\vTwelveStatus{C}{RETROSPECTIVE}{COMPATIBLE}
The 2026 validation matrix reports all nine displayed magnitudes within roughly
$1.71\sigma$ of its selected reference snapshot, while retaining postdiction
provenance.

\section{CKM CP phase and exact Jarlskog invariant}
The direct TIR phase assignment is
\begin{equation}
\boxed{
\delta_{\rm CKM}
=\arccos\frac{L_4}{L_5}
=\arccos\frac25
=66.4218215^\circ\ldots .
}
\label{eq:v12-ckm-delta}
\end{equation}

Once the three CKM mixing amplitudes and this phase are fixed, the Jarlskog
invariant is not an independent parameter.  Let
\begin{equation}
s_{12}=\lambda,
\qquad
s_{23}=|V_{cb}|,
\qquad
s_{13}=|V_{ub}|,
\qquad
c_{ij}=\sqrt{1-s_{ij}^2}.
\label{eq:v12-ckm-sines}
\end{equation}
For the standard three-angle one-phase CKM parameterization,
\begin{equation}
\boxed{
J_{\rm CP}^{\rm CKM}
=\operatorname{Im}
\left(V_{ud}V_{cs}V_{us}^{\ast}V_{cd}^{\ast}\right)
=s_{12}s_{23}s_{13}c_{12}c_{23}c_{13}^{2}\sin\delta_{\rm CKM}.
}
\label{eq:v12-ckm-jarlskog}
\end{equation}
Substituting Eqs.~\eqref{eq:v12-ckm-lambda}--\eqref{eq:v12-ckm-delta}
gives
\begin{equation}
\boxed{
J_{\rm CP}^{\rm CKM}
=2.97106576668\times10^{-5}.
}
\label{eq:v12-ckm-jarlskog-value}
\end{equation}
The dedicated closure validator constructs the full complex CKM matrix,
checks $VV^\dagger=I$ and $\det V=1$ to floating-point roundoff, and verifies
that all nine nontrivial rephasing-invariant quartets carry the same
$|J_{\rm CP}^{\rm CKM}|$ up to orientation sign.

\vTwelveStatus{C}{RETROSPECTIVE}{COMPATIBLE}
The external-data timing classification remains retrospective because the
underlying mixing-magnitude formulas retain their historical provenance.  The
Jarlskog reduction itself is an exact matrix identity conditional on those CKM
inputs.

\section{Historical independent Jarlskog assignment: deprecated}
Earlier TIR versions used the independent expression
\begin{equation}
J_{\rm old}
=\kappa^2\frac{L_4}{L_5}
\left(1-\frac{(L_4/L_5)^2}{2}\right)
=3.11011545905\times10^{-5}.
\label{eq:v12-ckm-j-old}
\end{equation}
It is not the exact Jarlskog invariant of the current CKM matrix and is no
longer retained as an independent CP assignment.  Its gap from the exact matrix
invariant is
\begin{equation}
\boxed{
\frac{|J_{\rm old}-J_{\rm CP}^{\rm CKM}|}{J_{\rm old}}
=4.47088522\%.
}
\label{eq:v12-ckm-j-old-gap}
\end{equation}

Writing
\[
c=\frac{L_4}{L_5}=\frac25,
\]
the exact phase factor is
\[
\sin\delta_{\rm CKM}=\sqrt{1-c^2},
\]
whereas $1-c^2/2$ is only the first nontrivial truncation of the Taylor series
of $\sqrt{1-c^2}$.  In addition, the historical assignment omitted the full
mixing Jacobian
\begin{equation}
\mathcal M
=\frac{s_{12}s_{23}s_{13}c_{12}c_{23}c_{13}^2}{\kappa^2c}
=0.95892344664\ldots .
\label{eq:v12-ckm-mixing-jacobian}
\end{equation}
Thus the exact current relation may be written as
\begin{equation}
\boxed{
J_{\rm CP}^{\rm CKM}
=\kappa^2c\,\mathcal M\sqrt{1-c^2}.
}
\label{eq:v12-ckm-j-factorized}
\end{equation}
The former independent assignment effectively replaced the product
$\mathcal M\sqrt{1-c^2}$ by $1-c^2/2$.

The archived Stage-32 phase backsolve also used a denominator proportional to
$s_{12}s_{23}s_{13}c_{12}c_{23}c_{13}$; the exact Jarlskog identity contains
$c_{13}^2$.  Because $c_{13}\simeq1$, the numerical effect is tiny, but the
formal correction is retained.

\section{Remaining CKM upstream gate}
The Jarlskog invariant is now closed downstream of the current CKM inputs.  A
full first-principles CKM closure still requires an independent forcing theorem
for the functional forms
\begin{equation}
\boxed{
\begin{aligned}
s_{12}
&=\frac{L_4}{L_3+L_4}+\frac{L_4}{L_3}\kappa,\\
s_{23}
&=\frac12\left(\frac{L_4}{L_3}\right)^2,\\
s_{13}
&=\frac12\left(\frac{L_4}{L_3}\right)^2
\frac{L_4}{L_3+L_4}\frac{L_4}{L_5}.
\end{aligned}
}
\label{eq:v12-ckm-upstream-open}
\end{equation}
The new invariant closure does not silently promote these three historical
structural assignments.

\section{GREMLIN Pass-4 selector audit}

The preserved 600-cell/McKay/A2 branch narrows the upstream problem without
changing its evidence class.  The exact mathematical layer now contains:

\begin{itemize}
\item tetrahedral incidence counts \(\nu_F=2\) and \(\nu_E=5\), with
      the structural ratio alphabet \((2/7,2/9,2/5)\);
\item a local centred \(A_2\) three-channel carrier and an exact
      non-adjacent-plane principal cosine \(2/9\);
\item exact finite-operator readouts for \(s_{23}=2/49\) and the
      composite-root factorization \(s_{13}=s_{23}(2/9)(2/5)\);
\item a conditional projective even-sector selector for the rational base
      weights;
\item a half-kernel boundary compression \(9\to7\) and a Kronecker-sum
      realization of the historical scalar readout \(2/9+(2/7)\kappa\);
\item an exact operator-identifiability no-go: the scalar trace alone does not
      determine a unique generator;
\item an exact angle--sine firewall: additivity of a Lie-algebra angle does not
      imply additivity of the sine coordinate.
\end{itemize}

These results are kept as structural/conditional mechanisms.  They do not turn
the historically developed CKM formula into a prospective prediction, and they
do not establish that the 600-cell is a literal physical lattice.  The
remaining physical theorem is a coefficient-free flavour/readout map that
selects the admitted channel operators and their observable sine coordinates
without consuming CKM targets.

\vTwelveStatus{B}{--}{OPEN}

The Pass-4 provenance chain is recorded in
\path{TIR/integration/TIR_CKM_600CELL_FLAG_INCIDENCE_CANDIDATE_V0_1.md}
through
\path{TIR/integration/TIR_CKM_SU3F_ANGLE_SINE_GENERATOR_FIREWALL_V0_8.md}.

\section{Reproducibility}
The retained arithmetic/internal-consistency audit is
\path{TIR/validation/tir_v12_flavour_mixing_internal_consistency_v0_1.py}.
The exact CKM/Jarlskog closure audit is
\path{TIR/validation/tir_ckm_jarlskog_invariant_closure_v0_1.py}.
The theorem surface is
\path{TIR/integration/TIR_CKM_JARLSKOG_INVARIANT_CLOSURE_V0_1.md}.
Current external comparison is read from the unified evidence layer rather than
from the historical tables embedded in the v11 sector chapters.

% ---- END TIR/monograph/v12/chapters/ch13_ckm_pmns_flavour_mixing.tex ----
\clearpage

% ---- BEGIN TIR/monograph/v12/chapters/ch14_hadrons_baryons_mesons.tex ----
% !TEX root = ../../tir_monograph_v12.tex
\chapter{Hadrons — Baryons and Mesons}
\label{ch:v12_hadrons}

\section{Source and status discipline}

This chapter consolidates the historical baryon and meson source family. Standard $SU(3)$ flavour relations are separated from TIR coefficient assignments and from empirical comparison. Current observable verdicts are owned by Chapter~\ref{ch:v12_evidence_matrix}.

Primary legacy sources are
\path{ch08_gmo_formula.tex} through \path{ch15_heavy_mesons.tex}. The v12 arithmetic audit is
\path{TIR/validation/tir_v12_hadron_formula_audit_v0_2.py}.

\section{Baryon octet: the standard GMO carrier}

For the baryon octet the standard Gell--Mann--Okubo organization is
\begin{equation}
M = a+bY+c\left[I(I+1)-\frac{Y^2}{4}\right],
\label{eq:v12_gmo_standard}
\end{equation}
with the familiar octet relation
\begin{equation}
\boxed{
\frac12(M_N+M_\Xi)=\frac34M_\Lambda+\frac14M_\Sigma.
}
\label{eq:v12_gmo_relation}
\end{equation}
\vTwelveStatus{A}{--}{PASS}
The algebraic relation belongs to the established $SU(3)$ flavour framework \cite{GellMann1962,okubo1962}.

The historical TIR parameterization writes
\begin{equation}
M=M_0+\alpha+\beta Y+\gamma\left[I(I+1)-\frac{Y^2}{4}\right].
\label{eq:v12_tir_gmo}
\end{equation}
Its coefficient family is retained as a retrospective model surface.

\section{Octet reference scale: v12 derivation audit}

The legacy source defines
\begin{equation}
S_0=\frac{s+u}{L_3}\,\kappa
\label{eq:v12_legacy_s0}
\end{equation}
and cites the exponential architecture
\begin{equation}
M_0=E_P\exp(-S_0/\kappa).
\label{eq:v12_legacy_m0_exp}
\end{equation}
Equations~\eqref{eq:v12_legacy_s0}--\eqref{eq:v12_legacy_m0_exp} imply
\begin{equation}
\frac{S_0}{\kappa}=\frac{s+u}{L_3}=\frac{10}{7},
\end{equation}
so literal evaluation gives
\begin{equation}
E_Pe^{-10/7}\simeq2.926\times10^{21}\ \mathrm{MeV}.
\end{equation}

The printed $925.95$ MeV value instead comes from a distinct proton-anchored linear map,
\begin{equation}
\boxed{
M_0^{\rm lin}
=E_{\rm proton}
\left(1-\frac{s+u}{L_3}\kappa\right)
\simeq925.9496\ \mathrm{MeV}.
}
\label{eq:v12_octet_m0_linear}
\end{equation}
The corresponding exact proton-anchored exponential would be
\begin{equation}
E_{\rm proton}\exp\!\left[-\frac{s+u}{L_3}\kappa\right]
\simeq926.0302\ \mathrm{MeV}.
\end{equation}

The v12 audit therefore preserves Eq.~\eqref{eq:v12_octet_m0_linear} as the formula that generated the historical table while routing its derivational parentage to repair.
\vTwelveStatus{D}{RETROSPECTIVE}{QUARANTINED}

\section{Octet coefficient family}

The historical coefficient assignments are
\begin{align}
\alpha
&=E_{\rm proton}\kappa
\left(q_s^2-q_u^2-q_d^2+L_3\right),
\label{eq:v12_octet_alpha}\\
\beta
&=-\alpha\frac{L_3^2-1}{L_3^2},
\label{eq:v12_octet_beta}\\
\gamma
&=\alpha\frac{L_4(L_3-L_4)}{L_3^2}.
\label{eq:v12_octet_gamma}
\end{align}
For the canonical labels these give approximately
\[
\alpha=189.765\ \mathrm{MeV},\qquad
\beta=-185.892\ \mathrm{MeV},\qquad
\gamma=38.728\ \mathrm{MeV}.
\]

The legacy chapter attributes Eqs.~\eqref{eq:v12_octet_beta}--\eqref{eq:v12_octet_gamma} to a ``J--KJ'' eigenvalue construction. Version 12 records a completion gate for the explicit operator/matrix, its domain and its eigenvalue calculation. Until that surface is supplied, the coefficient equations are typed retrospective assignments rather than an operator-level theorem.
\vTwelveStatus{C}{RETROSPECTIVE}{QUARANTINED}

With Eqs.~\eqref{eq:v12_octet_m0_linear}--\eqref{eq:v12_octet_gamma}, the constructed masses satisfy Eq.~\eqref{eq:v12_gmo_relation} identically. This is an internal algebraic consistency check of the chosen GMO parameterization. Empirical independence is controlled separately through provenance.

\section{Baryon decuplet}

The historical decuplet uses
\begin{equation}
M=M_0'+\alpha_{10}+\beta_{10}Y,
\label{eq:v12_decuplet}
\end{equation}
with
\begin{align}
M_0'&=E_{\rm proton}\left[1+(s-u)\kappa\right],\\
\alpha_{10}
&=E_{\rm proton}\kappa
\frac{L_3^2+L_4^2+L_5^2+L_3+L_4+L_5+L_4}{L_4},\\
\beta_{10}
&=-E_{\rm proton}\kappa\frac{L_3L_5}{L_4}.
\end{align}
The displayed values are approximately
\[
M_0'=972.72\ \mathrm{MeV},\qquad
\alpha_{10}=405.41\ \mathrm{MeV},\qquad
\beta_{10}=-150.95\ \mathrm{MeV}.
\]
They generate the equal-spaced sequence reported historically for $\Delta$, $\Sigma^*$, $\Xi^*$ and $\Omega$.

The proton mass enters explicitly as an external anchor and the coefficient numerators are discrete structural assignments. The frozen 2026 evidence matrix therefore treats this as postdiction with a proton anchor and candidate compatibility.
\vTwelveStatus{C}{RETROSPECTIVE}{COMPATIBLE}

\section{Pseudoscalar mesons: retained arithmetic failure}

The legacy pion and kaon source prints
\begin{align}
\frac{\Delta S_\pi}{\kappa}&=\frac6\pi,
& m_\pi&=E_Pe^{-6/\pi},\\
\frac{\Delta S_K}{\kappa}&=\zeta(2)-1,
& m_K&=E_Pe^{-(\zeta(2)-1)}.
\end{align}
Literal evaluation with the displayed $E_P=1.22089\times10^{22}$ MeV gives
\begin{align}
E_Pe^{-6/\pi}&\simeq1.808\times10^{21}\ \mathrm{MeV},\\
E_Pe^{-(\zeta(2)-1)}&\simeq6.406\times10^{21}\ \mathrm{MeV}.
\end{align}
These values establish the formula/value failure recorded by the PDG-2026 audit.
\vTwelveStatus{D}{RETROSPECTIVE}{FAIL}

The $\eta$--$\eta'$ source introduces the mixing matrix
\begin{equation}
\mathcal M^2=
\begin{pmatrix}
m_8^2&M_{81}^2\\
M_{81}^2&m_1^2
\end{pmatrix},
\end{equation}
while the numerical entries and their independent derivation remain the active provenance gate for the displayed eigenvalues.
\vTwelveStatus{C}{RETROSPECTIVE}{QUARANTINED}

\section{Heavy and vector mesons}

The historical heavy/vector chapter lists near-identical ``predicted'' and PDG masses for open- and hidden-flavour states while referring to a zeta-operator spectrum. Version 12 requires the operator, input map and reproducible calculation that generates those values before the table can carry evidential weight.
\vTwelveStatus{C}{RETROSPECTIVE}{QUARANTINED}

\section{Hadronic evidence boundary}

The v12 hadronic hierarchy is therefore
\[
\boxed{\begin{aligned}
\text{standard }SU(3)\text{ mass algebra}
&\to \text{TIR coefficient/anchor assignments}\\
&\to \text{reproducible mass calculation}\\
&\to \text{Chapter 19 verdict}.
\end{aligned}}
\]

The audit receipt is
\path{TIR/validation/TIR_V12_HADRON_FORMULA_AUDIT_V0_2.json}. It records \texttt{PASS\_WITH\_QUARANTINES}: the diagnostic checks themselves pass, while the identified legacy formula/provenance surfaces retain their explicit v12 quarantine/failure states. Version 0.2 supersedes the v0.1 validator metric while preserving the same audited scientific witnesses.

% ---- END TIR/monograph/v12/chapters/ch14_hadrons_baryons_mesons.tex ----
\clearpage

% ---- BEGIN TIR/monograph/v12/chapters/ch15_gauge_hypercharge_anomalies.tex ----
% !TEX root = ../../tir_monograph_v12.tex
\chapter[Gauge Structure and Anomalies]{Gauge Structure, Hypercharge and Anomaly Cancellation}
\label{ch:v12_gauge_hypercharge}

\section{Typed gauge transport parent}

The current TIR transport layer uses the common connection-holonomy family
\[
W_{ij}^{(G,R)}
=\mathcal P\exp\!\left(\int_{\gamma_{ij}}A_R\right),
\]
with sector-typed instances including
\[
W_{ij}^{WT}\in U(1),
\qquad
W_{ij}^{X}\in SU(2),
\qquad
W_{ij}^{c}\in SU(3).
\]
The source binding and transformation laws are recorded in
\path{TIR/foundations/TIR_WIJ_HOLONOMY_CROSSWALK_V0_1.md}.

This chapter audits the separate representation/hypercharge layer used by the Standard Model branch.

\section{Hypercharge assignment map}

The legacy TIR formulas are
\begin{align}
Y_Q&=\frac{q_s}{L_3L_4N_c}=\frac16,\\
Y_{uR}&=\frac{L_4}{N_c}=\frac23,\\
Y_{dR}&=-\frac{L_4}{N_cL_4}=-\frac13,\\
Y_{eR}&=-\frac{L_3-L_4}{L_5}=-1,\\
Y_L&=-\frac12\qquad\text{(external input in the legacy surface)}.
\end{align}
For
\[
(q_s,L_3,L_4,L_5,N_c)=(7,7,2,5,3)
\]
the arithmetic map reproduces the conventional one-generation Standard Model hypercharges exactly.

\vTwelveStatus{B}{--}{OPEN}
The arithmetic map is reproducible. Version 12 keeps the derivation/uniqueness gate explicit because $Y_L$ enters as an external parent and several formulas simplify by cancellation of the same structural labels appearing in numerator and denominator.

\section{Anomaly convention}

For anomaly sums, all fermions are represented as left-handed Weyl fields. A right-handed field of hypercharge $Y_R$ is therefore represented by its left-handed charge conjugate with hypercharge $-Y_R$. With this convention the four local one-generation anomaly coefficients are
\begin{align}
\mathcal A_{33Y}
&=2Y_Q-Y_{uR}-Y_{dR},\\
\mathcal A_{22Y}
&=N_cY_Q+Y_L,\\
\mathcal A_{YYY}
&=N_c\left(2Y_Q^3-Y_{uR}^3-Y_{dR}^3\right)
 +\left(2Y_L^3-Y_{eR}^3\right),\\
\mathcal A_{\mathrm{grav}\,Y}
&=N_c\left(2Y_Q-Y_{uR}-Y_{dR}\right)
 +\left(2Y_L-Y_{eR}\right).
\end{align}

Substituting
\[
Y_Q=\frac16,\quad
Y_{uR}=\frac23,\quad
Y_{dR}=-\frac13,\quad
Y_L=-\frac12,\quad
Y_{eR}=-1
\]
gives
\begin{align}
\mathcal A_{33Y}&=0,\\
\mathcal A_{22Y}&=0,\\
\mathcal A_{YYY}&=0,\\
\mathcal A_{\mathrm{grav}\,Y}&=0.
\end{align}

\vTwelveStatus{A}{--}{PASS}
These cancellations are the standard anomaly consistency of the displayed Standard Model representation content. Their exact rational audit is reproduced by
\path{TIR/validation/tir_v12_hypercharge_anomaly_audit_v0_1.py}.

\section[Global SU(2) consistency]{Global $SU(2)$ consistency}

For one generation, the left-handed $SU(2)$ doublets consist of three colour copies of $Q_L$ plus one lepton doublet $L_L$:
\[
N_{\rm doublet}=N_c+1=4.
\]
The count is even, satisfying the Witten global $SU(2)$ consistency condition.
\vTwelveStatus{A}{--}{PASS}

\section{Baryon-current anomaly and legacy A5 repair}

The legacy Chapter 28 labels
\begin{equation}
3Y_Q-Y_{uR}-Y_{dR}=\frac16
\label{eq:v12_legacy_a5}
\end{equation}
as a fifth condition $[SU(2)]^2U(1)_B$. Equation~\eqref{eq:v12_legacy_a5} is built from hypercharges. The baryon-current coefficient instead uses baryon charge $B=1/3$ for each coloured quark doublet. In the conventional $T(\mathbf2)=1/2$ normalization,
\begin{equation}
\boxed{
\mathcal A_{SU(2)^2B}
=N_c\left(\frac13\right)T(\mathbf2)
=3\left(\frac13\right)\left(\frac12\right)
=\frac12.
}
\end{equation}

The nonzero baryon-current anomaly is the expected electroweak anomalous current structure; the v12 repair is the charge/type correction of the legacy A5 expression.
\vTwelveStatus{D}{RETROSPECTIVE}{QUARANTINED}

\section{Result hierarchy}

The v12 gauge/anomaly dependency is
\[
\boxed{\begin{aligned}
\text{typed }W_{ij}\text{ transport}
&\to \text{representation + hypercharge assignment}\\
&\to \text{left-Weyl anomaly sums}\\
&\to \text{current evidence/status layer}.
\end{aligned}}
\]

The exact standard anomaly cancellations and Witten count receive formal \texttt{PASS}. The TIR hypercharge source map remains a structural \texttt{OPEN} gate for derivation/uniqueness, and the legacy A5 label is quarantined with an explicit corrected baryon-current coefficient.

Audit receipt:
\path{TIR/validation/TIR_V12_HYPERCHARGE_ANOMALY_AUDIT_V0_1.json}.

% ---- END TIR/monograph/v12/chapters/ch15_gauge_hypercharge_anomalies.tex ----
\clearpage

% ---- BEGIN TIR/monograph/v12/chapters/ch16_electroweak_higgs.tex ----
% !TEX root = ../../tir_monograph_v12.tex
\chapter{Electroweak and Higgs Sector}
\label{ch:v12_electroweak_higgs}

\section{Structural scale and coupling candidates}

The historical electroweak branch supplies the structural scale
\begin{equation}
v_0
=E_{\rm proton}\left(L_3^2L_5+L_3L_4+L_4\right).
\label{eq:v12_ew_v0}
\end{equation}
For $(L_3,L_4,L_5)=(7,2,5)$ the integer factor is $261$, giving
\begin{equation}
\boxed{v_0=244.888992\ \mathrm{GeV}}.
\end{equation}
The proton mass is an explicit external anchor in this relation.
\vTwelveStatus{B}{RETROSPECTIVE}{OPEN}

The weak-angle candidate is
\begin{equation}
\boxed{
\sin^2\theta_W^{(0)}
=\frac{L_4}{L_3+L_4}+\kappa
=\frac29+\kappa
=0.2314153722\ldots .
}
\label{eq:v12_sin2w0}
\end{equation}
Its precision comparison requires a declared renormalization scheme and scale.
\vTwelveStatus{B}{RETROSPECTIVE}{TENSION}

The separate $SU(2)$ coupling candidate is
\begin{equation}
\boxed{
g_0=\frac{L_4}{L_3}+\frac{L_4}{L_5}
=\frac{24}{35}
=0.6857142857\ldots .
}
\label{eq:v12_g0}
\end{equation}

\section[Tree-level W/Z map]{Tree-level $W/Z$ map}

Using Eqs.~\eqref{eq:v12_ew_v0}--\eqref{eq:v12_g0}, the historical tree-level construction is
\begin{align}
M_W^{(0)}&=\frac{g_0v_0}{2}
=83.9619401\ldots\ \mathrm{GeV},\\
M_Z^{(0)}&=\frac{M_W^{(0)}}{\sqrt{1-\sin^2\theta_W^{(0)}}}
=95.7715725\ldots\ \mathrm{GeV}.
\end{align}
The frozen PDG-2026 comparison assigns precision verdict \texttt{FAIL} to both pole-mass readings.
\vTwelveStatus{C}{RETROSPECTIVE}{FAIL}

The active closure target is the joint transport
\begin{equation}
\boxed{
(g_0,\theta_W^{(0)},v_0)
\xrightarrow{\mathcal R_{EW}(\mu,\mathrm{scheme})}
(g(\mu),\theta_W(\mu),v(\mu))
\longrightarrow
(M_W^{\rm pole},M_Z^{\rm pole}).
}
\label{eq:v12_rew}
\end{equation}
This transport is a shared sector gate rather than an observable-specific correction.

\section{Fine-structure relation}

The independent TIR relation is
\begin{equation}
\boxed{
\alpha_{\rm TIR}^{-1}
=(L_3L_4)^2-L_3^2-L_4L_5+L_4^2\kappa
=137.0367726000\ldots .
}
\label{eq:v12_alpha_inv}
\end{equation}
Its interpretation as the zero-momentum fine-structure constant carries the frozen precision verdict \texttt{FAIL}.
\vTwelveStatus{C}{RETROSPECTIVE}{FAIL}

\section{Electroweak charge-closure gate}

At tree level the standard electroweak relation between the electromagnetic and $SU(2)$ couplings is
\begin{equation}
e=g\sin\theta_W,
\qquad
\alpha=\frac{e^2}{4\pi}.
\label{eq:v12_ew_charge_relation}
\end{equation}
Applying the raw structural pair $(g_0,\theta_W^{(0)})$ gives
\begin{equation}
e_{g\theta}=g_0\sin\theta_W^{(0)}
=0.3298673257\ldots,
\end{equation}
and therefore
\begin{equation}
\boxed{
\alpha_{g\theta}^{-1}
=\frac{4\pi}{e_{g\theta}^2}
=115.4865120\ldots .
}
\label{eq:v12_alpha_from_gtheta}
\end{equation}
By contrast, Eq.~\eqref{eq:v12_alpha_inv} gives $137.0367726\ldots$. Equivalently, combining the TIR $\alpha$ relation with Eq.~\eqref{eq:v12_sin2w0} requires
\begin{equation}
\boxed{g_{\alpha\theta}=0.6294920778\ldots}
\end{equation}
rather than the raw $g_0=24/35$.

The three raw structural relations therefore define distinct pre-transport quantities. Their common electroweak binding is an explicit completion problem:
\begin{equation}
\boxed{
(g_0,\theta_W^{(0)},\alpha_0,v_0)
\xrightarrow{\mathcal R_{EW}(\mu,\mathrm{scheme})}
(g,e,\theta_W,\alpha,v)_{\mu,\mathrm{scheme}}
}
\label{eq:v12_ew_full_closure}
\end{equation}
subject to Eq.~\eqref{eq:v12_ew_charge_relation} at the declared comparison surface.
\vTwelveStatus{B}{--}{OPEN}

\section{Higgs mass relation}

The current Higgs assignment is
\begin{equation}
\boxed{
M_H=v_0\kappa(L_3^2+L_4+L_5)
=v_0\kappa\,56
=126.0728693\ldots\ \mathrm{GeV}.
}
\label{eq:v12_higgs_mass}
\end{equation}
The functional form was introduced after the earlier scalar relation had failed and after the Higgs target was known, so its evidence timing is retrospective. The frozen 2026 precision comparison assigns verdict \texttt{FAIL}.
\vTwelveStatus{C}{RETROSPECTIVE}{FAIL}

The structural completion target is a scalar-sector derivation of the discrete factor $L_3^2+L_4+L_5$ from the field/action geometry, followed by the same scheme/scale discipline used by the wider electroweak sector.

\section{v12 electroweak hierarchy}

The resulting dependency structure is
\[
\boxed{
\begin{array}{c}
(v_0,g_0,\theta_W^{(0)},\alpha_0)_{\rm structural}\\[2pt]
\downarrow\\[-2pt]
\mathcal R_{EW}(\mu,\mathrm{scheme})\\[2pt]
\downarrow\\[-2pt]
(g,e,\theta_W,\alpha,v)_{\mu,\mathrm{scheme}}\\[2pt]
\downarrow\\[-2pt]
(M_W^{\rm pole},M_Z^{\rm pole},M_H^{\rm observable})
\end{array}}
\]
with the current empirical verdicts maintained in Chapter~\ref{ch:v12_evidence_matrix}.

The audit receipt is
\path{TIR/validation/TIR_V12_ELECTROWEAK_HIGGS_AUDIT_V0_1.json}. It records \texttt{PASS\_WITH\_EW\_CLOSURE\_GATE}: the arithmetic reproduction checks pass and the raw electroweak charge-binding mismatch is promoted to an explicit completion gate.

% ---- END TIR/monograph/v12/chapters/ch16_electroweak_higgs.tex ----
\clearpage

% ---- BEGIN TIR/monograph/v12/chapters/ch17_strong_cp_neutron_edm.tex ----
% !TEX root = ../../tir_monograph_v12.tex
\chapter{Strong CP and Neutron EDM}
\label{ch:v12_strong_cp_edm}

\section{Frozen strong-CP map}

The reviewed v11 branch supplies the assignment
\begin{equation}
\boxed{
\theta_{\rm QCD}
=\kappa\left(\frac{L_4}{L_3}\right)^{L_3+L_4+L_5}
=\kappa\left(\frac27\right)^{14}.
}
\label{eq:v12_theta_qcd}
\end{equation}
The exponent $14=L_3+L_4+L_5$ and the suppression operator define the frozen retrospective strong-CP version.
\vTwelveStatus{C}{RETROSPECTIVE}{OPEN}

The current source-completion target is supplied by the non-Abelian colour-holonomy branch:
\begin{equation}
\boxed{
W_{ij}^{c}
\longrightarrow
U_\gamma
\longrightarrow
\text{topological CP phase / sector invariant}
\longrightarrow
\theta_{\rm QCD}.
}
\label{eq:v12_strongcp_source_axis}
\end{equation}
This gives the existing numerical assignment a precise upstream theorem target.

\section{Corrected arithmetic}

The exact numerical suppression in Eq.~\eqref{eq:v12_theta_qcd} is
\begin{equation}
\left(\frac27\right)^{14}
=2.4157243620710218\times10^{-8}.
\end{equation}
With
\[
\kappa=\frac{\ln2}{24\pi}=0.009193150006360484\ldots,
\]
the frozen phase is
\begin{equation}
\boxed{
\theta_{\rm QCD}
=2.220811643453839\times10^{-10}.
}
\label{eq:v12_theta_qcd_numeric}
\end{equation}
The v12 validator reproduces this value independently from the published decimal.
\vTwelveStatus{C}{RETROSPECTIVE}{PASS}
Here \texttt{PASS} is the arithmetic/formula evaluation of the frozen assignment; its physical consequence is tested separately below.

\section{Fixed hadronic conversion}

The publication snapshot uses the conversion coefficient
\begin{equation}
C_n=2.4\times10^{-16}\ e\,\mathrm{cm},
\end{equation}
with
\begin{equation}
\boxed{
d_n=C_n\theta_{\rm QCD}.
}
\label{eq:v12_nedm_map}
\end{equation}
Therefore
\begin{equation}
\boxed{
d_n^{\rm frozen}
=5.329947944289213\times10^{-26}\ e\,\mathrm{cm}.
}
\label{eq:v12_nedm_value}
\end{equation}
The coefficient is an explicit external/hadronic matching input to this version.

\section{Physical gate}

Against the manuscript bound
\begin{equation}
|d_n|<1.8\times10^{-26}\ e\,\mathrm{cm}
\qquad(90\%\ \mathrm{CL}),
\label{eq:v12_nedm_bound}
\end{equation}
the ratio is
\begin{equation}
\boxed{
\frac{d_n^{\rm frozen}}{d_n^{\rm bound}}
=2.96108219127\ldots .
}
\end{equation}
The frozen physical gate therefore has verdict
\vTwelveStatus{D}{RETROSPECTIVE}{FAIL}
The failure remains attached to the exact exponent, normalization and hadronic conversion that generated it.

\section{Versioned replacement rule}

A future strong-CP relation is promoted as a new version only after its upstream source map is frozen. The preferred structural route is Eq.~\eqref{eq:v12_strongcp_source_axis}, followed by a declared hadronic conversion and the same observable-level nEDM gate. This produces an append-only chain
\[
\boxed{
\text{colour holonomy source}
\to\theta_{\rm QCD}
\to\text{hadronic matching}
\to d_n
\to\text{experimental verdict}.
}
\]

\section{Audit receipt}

The arithmetic and empirical gate are reproduced by
\path{TIR/validation/tir_v12_strong_cp_nedm_audit_v0_1.py}.
Its receipt
\path{TIR/validation/TIR_V12_STRONG_CP_NEDM_AUDIT_V0_1.json}
reports \texttt{PASS\_ARITHMETIC\_PHYSICAL\_FAIL\_RETAINED}: the frozen calculation is reproducible and the physical \texttt{FAIL} remains visible in Chapter~\ref{ch:v12_evidence_matrix}.

% ---- END TIR/monograph/v12/chapters/ch17_strong_cp_neutron_edm.tex ----
\clearpage
